\section{A 62-operation quadratic system and its quartic}
\label{sec:quartic}

For a presentation with small degree, the most readable new result is
a system of 22 quadratic equations in 31 strictly positive integers.
Three coordinated changes of coordinates reduce package 30's system
from 65 operations to 62. Its sum of squares is a single quartic using
127 operations, compared with 130 before the changes.

The 31 variables are
\begin{align*}
 &J,F,\alpha,z,f,h,i,o,j,s,w,\tau,\eta,\zeta,y,\gamma,\\
 &N,E,T,c,a,X,D,C_2,S,F_2,V,g,I,B_a,\Delta.
\end{align*}
All are strictly positive. Use the four fixed coefficients
\eqref{eq:new-ports} and the affine abbreviations
\[
 q=\beta J+1,\qquad k=\eta+\zeta,\qquad H=4a+3,
 \qquad C=q-F-\alpha.
\]
The equations are
\begin{align}
 S&=q^2, &N&=qS, &X&=wq,\label{eq:quad-start}\\
 a&=E+sN, &E&=X(a-E), &c&=k(a-E)+\eta,\\
 T&=E(c-\eta), &\tau^2&=T(T+k)+1,\\
 D&=X+ac+\gamma H, &\Delta&=a^2+H,\\
 C_2&=c^2, &D^2&=\Delta C_2+1,\\
 k&=jc-V+1+hE.&&
\end{align}
The two packing and transport equations are
\begin{align}
 jc-V&=(q-F)(N-q)-C(S-1)+(M_C+qM_F)J,\\
 (K_0+w)C&=F+z(q-1).
\end{align}
The remaining seven equations are
\begin{align}
 F_2&=f^2, &V&=of-c, &g&=y^2-V^2,\\
 I&=iC_2, &B_a&=I^2, &B_a&=\Delta(F_2-1),\\
 y^2-B_ag&=1.&&&\label{eq:quad-end}
\end{align}
Since every abbreviation is affine and every displayed product has
at most two affine factors, every residual has total degree at most two.
Writing the 22 residuals as $r_j$, the single polynomial
\begin{equation}\label{eq:new-quartic}
 Q_{31}=\sum_{j=1}^{22}r_j^2
\end{equation}
has degree exactly four: the coefficient of $\tau^4$ is one. Over the
stated positive domain, $Q_{31}=0$ is equivalent to all 22 equations.

\subsection{The three shifts and the positive zero sets}

The predecessor coordinates being changed are denoted by
$L_{\rm old},P_{\rm old},z_{\rm old}$. The exact substitutions are
\[
 L_{\rm old}=S-1,\qquad P_{\rm old}=T+k,
 \qquad z_{\rm old}=z+1.
\]
All other variables are untouched.

The old equations for $L_{\rm old}$ and $N$ become $S=q^2$ and
$N=qS$. In the packing equation $L_{\rm old}$ becomes $S-1$.
This removes two gates and introduces one, for a net saving of one.
The old first-root definition was
$P_{\rm old}=E(c-\eta)+k$; after the shift it is
$T=E(c-\eta)$. The first norm becomes $\tau^2=T(T+k)+1$.
This saves a second gate. Finally the transport identity
\[
 (K_0+w)C+(q-F)-(z+1)(q-1)-1
       =(K_0+w)C-F-z(q-1)
\]
saves a third gate when its two sides are compared directly.
These are polynomial identities. The sparse-ring audit verifies each of
the 22 residuals against both the predecessor's expression and its
literal sparse coefficient data after substitution.

Given a new positive solution, $S=q^2>1$ implies
$L_{\rm old}>0$, while $P_{\rm old}=T+k>0$ and
$z_{\rm old}=z+1>0$. The identities give a positive predecessor solution.
Conversely, on a positive predecessor solution,
$S=L_{\rm old}+1=q^2$, and
\[
 a-E=sN>0,\qquad c-\eta=k(a-E)>0,
 \qquad T=P_{\rm old}-k=E(c-\eta)>0.
\]
For the final shift, the inherited Pell/compiler theorem supplies
$C>0$, $X=wq=2^R$, $R>3q$, and $F<q$.
Consequently $w>q>F$, and the transport equation gives
\[
 (z_{\rm old}-1)(q-1)=(K_0+w)C-F>0.
\]
Thus $z_{\rm old}\ge2$, so the new $z$ is positive.
This positivity step is a use of the inherited arithmetic theorem,
not an inference from a finite test of a few tuples.
The transformations are inverse on the positive zero sets.

\begin{theorem}[The low-degree construction]
ZFC proves that the 22-equation system above has no positive integer
solution if and only if $\Con(\ZFC)$. Equivalently, ZFC proves the
same statement for the quartic $Q_{31}=0$.
\end{theorem}
\begin{proof}
The shifts give the proved bijection to package 30's positive quartic
system. The complete coefficient replacement theorem permits the new
tuple in that system. Its inherited Pell/compiler equivalence and
the literal K215 calibration give the theorem.
\end{proof}

\subsection{Cost and the quartic used by the Jones branch}

The system has 62 live binary gates: 33 multiplications and 29 additions
or subtractions. The sum-of-squares finalizer adds 22 residual
subtractions, 22 squares, and 21 additions, giving 127 gates.
Its expansion has 297 monomials when the four coefficient ports remain
independent formal symbols; the predecessor has 306 under that convention.
This count does not identify algebraic coefficient ports with quantified
unknowns.

Replacing every positive variable by its nonnegative successor gives a
normalized quartic: at all positive variables equal to one the residual
$E-X(a-E)$ equals one independently of the coefficient tuple. Thus the
shifted quartic is nonzero at the all-zero tuple, as required by Jones.
To keep the index coefficients and their exact source map simple, the
signed construction below uses the predecessor's literal 31-variable
quartic, rebound to the same new tuple. It has the same existential
meaning. Its separate source data are included, and are not confused
with the new 297-term quartic.
